\begin{titlepage}
    \begin{center}
        \zihao{-1} \bfseries
        \titleEn \\
        \vspace*{36pt}
        \setmainfont{Times New Roman}
        \zihao{3}
        \begin{tabular}{ll}
            \makebox[6em][l]{Major: } &
            \makebox[10em][l]{Communication Studies} \\
            \makebox[6em][l]{Candidate: } &
            \makebox[10em][l]{\student} \\
            \makebox[6em][l]{Supervisor: } &
            \makebox[10em][l]{\teacher, Professor} \\
        \end{tabular}
        \vspace*{26pt}
    \end{center}

    \section*{Abstract}
    \phantomsection
    \addcontentsline{toc}{section}{Abstract}

    \subsection*{Background}
    The migration from traditional media to platform-based communication has profoundly changed the career trajectories of East Asian female celebrities. Visibility is no longer secured by print, drama, and variety programmes alone; instead, social media and creator-led video platforms reconfigure exposure, feedback, and crisis diffusion. Public cases therefore provide a useful environment for modelling career stages and resilience without relying on private data.

    \subsection*{Objective}
    \begin{enumerate}[label=\textbf{\arabic*.}]
        \item To build a reusable cross-media career life-cycle model;
        \item To summarize the mechanisms shaping successful and unstable transition;
        \item To demonstrate an interpretable prediction framework with fully public content.
    \end{enumerate}

    \subsection*{Methods}
    Public reports, platform timelines, programme archives, and interviews were reorganized into a four-layer workflow of stage segmentation, event coding, capability mapping, and resilience scoring. Descriptive comparison, rule-based scoring, and structured exhibits were used to convert qualitative observations into a thesis-ready analytical narrative.

    \subsection*{Results}
    The demonstration case showed a clear stage structure: an early visibility-building phase in traditional media, an expansion phase with diversified screen roles, a crisis phase characterized by declining narrative control, and a restart phase driven by self-managed media channels. Platform autonomy, role diversity, and output continuity emerged as three decisive factors for stable transition. A simple interpretable predictor based on these dimensions was sufficient to separate stable expansion, temporary fluctuation, and post-crisis restart trajectories.

    \subsection*{Conclusion}
    A cross-media life-cycle model helps transform complex celebrity narratives into comparable analytical stages, while an interpretable predictor offers a practical bridge from description to structured foresight. For this repository, the public example also demonstrates that a full doctoral-thesis layout can be showcased without exposing any real unpublished manuscript.

    \subsection*{Keywords}
    cross-media transition; career life cycle; image resilience; public case study; interpretable prediction

\clearpage{\cleardoublepage}
\end{titlepage}
